\chapter{Arbitrage theory}\label{sec:arbitrage}

We can generalize the approach used to derive the Black-Scholes PDE for a more general products
Thus in this section will try to give a general framework to derive equations to price a \textbf{contingent claim}.
A contingent claim is a financial instrument whose payoff depends on the outcome of uncertain future events,
for example an European option.

Throughout the chapter we work on a filtered probability space $(\Omega, \F, \{ \F_t \}_{t \in [0, T]}, \mathbb{P})$,
where $\{ \F_t \}$ is the filtration generated by a $k$-dimensional Brownian motion $W(t)$,
so $\F_t$ is the information available from the market up to time $t$.

\begin{definition}[Market]\label{def:market}
    A market is a set of $n + 1$ assets $S_0, S_1, \dots, S_n$,
    where the price of the riskless $S_0$ is given by
    \begin{equation*}
        \frac{dS_0(t)}{S_0} = r(S_0(t), t) dt,
    \end{equation*}
    and the prices of the other risky assets are given by
    \begin{equation*}
        dS_i(t) = \mu_i(S(t), t) dt + \sigma_i(S(t), t)^\intercal dW(t),
        \quad i = 1, \dots, n,
    \end{equation*}
    where $\mu_i$ is real valued and $\sigma_i$ takes values in $\R^k$.
\end{definition}

\begin{definition}[Portfolio]\label{def:portoflio}
    A portfolio in the market $\{ S_i(t) \}$ is an $n + 1$ dimensional $\F_t$-adapted stochastic process
    \begin{equation*}
        \theta(t) = (\theta_0(t), \theta_1(t), \dots, \theta_n(t)),
    \end{equation*}
    and its value at time $t$ is given by
    \begin{equation*}
        V(t) = \theta(t)^\intercal S(t)
            = \sum_{i = 0}^n \theta_i(t) S_i(t).
    \end{equation*}
\end{definition}

Note that as a portfolio $\theta(t)$ is an adapted process it doesn't use future information,
so it is  a generalization on what in reality would be to rebalance ones investments in the market at time $t$.

We can generalize the self-financing condition presented in \eqref{eq:self_financing} as,
\begin{equation}\label{eq:self_financing_portfolio}
    V(t) = V(0) + \int_0^t \theta(u)^\intercal dS(u),
\end{equation}
where the integral is a vector of stochastic integrals taken component wise,
\begin{equation*}
    \int_0^t \theta(u)^\intercal dS(u) = \sum_{i = 0}^n \int_0^t \theta_i(u) dS_i(u),
\end{equation*}
in differential form
\begin{equation}\label{eq:self_financing_portfolio_diff}
    dV(t) = \theta(t)^\intercal dS(t).
\end{equation}

As our market consists of assets that are described as stochastic processes,
an European contingent claim paying at time $T$ is a function $f(S(T), T)$ depending on these assets,
which is a random variable as $S(T)$ is a random vector of the values of the assets at time $T$. 

\section{PDE method}\label{subsec:pde_method}

Consider an European contingent claim with payoff $f(S(T))$,
similar to what we did in section \ref{subsec:black_scholes_model},
we want to construct a replicating portfolio $V(S(t), t)$ that has the same payoff as the contingent claim at time $T$.
We follow \cite[Ch.~1]{glasserman_monte_carlo}, see also \cite[Ch.~13]{bjork_arbitrage} for a detailed treatment.
Thus using the Ito formula \eqref{eq:ndim_ito_formula} we have
\begin{equation*}
\begin{aligned}
    dV(S(t), t)
        & = \frac{\partial V}{\partial t} dt
            + \sum_{0 \leq i \leq n} \frac{\partial V}{\partial S_i} dS_i(t)
            + \frac{1}{2}  \sum_{0 \leq i,j \leq n} \frac{\partial^2 V}{\partial S_i \partial S_j} dS_i(t) dS_j(t) \\
        & = \sum_{0 \leq i \leq n} \frac{\partial V}{\partial S_i} dS_i(t)
            + \left( \frac{\partial V}{\partial t} + \frac{1}{2}  \sum_{0 \leq i,j \leq n} \frac{\partial^2 V}{\partial S_i \partial S_j} \sigma_i(S(t), t)^\intercal \sigma_j(S(t), t) \right) dt,
\end{aligned}
\end{equation*}
where $\sigma_0 \equiv 0$ as $S_0$ is riskless.
Imposing the self-financing condition \eqref{eq:self_financing_portfolio_diff} we find that
\begin{equation}\label{eq:replicating_portfolio_pde_1}
    \theta_i(t) = \frac{\partial V}{\partial S_i}, \quad \forall i = 0, \dots, n,
\end{equation}
and 
\begin{equation}\label{eq:replicating_portfolio_pde_2}
    \frac{\partial V}{\partial t} + \frac{1}{2}  \sum_{0 \leq i,j \leq n} \frac{\partial^2 V}{\partial S_i \partial S_j} \sigma_i(S(t), t)^\intercal \sigma_j(S(t), t) = 0.
\end{equation}
Also using $V(S(t), t) = \theta(t)^\intercal S(t)$ in \eqref{eq:replicating_portfolio_pde_1} we get
\begin{equation}\label{eq:replicating_portfolio_pde_3}
    V(S(t), t) = \sum_{i = 0}^n \frac{\partial V}{\partial S_i} S_i.
\end{equation}
The equations \eqref{eq:replicating_portfolio_pde_2} and \eqref{eq:replicating_portfolio_pde_3} describe the PDE for the price of the contingent claim,
with independent variables $(S, t) \in [0, \infty)^{n + 1} \times [0, T]$ and terminal condition the payoff at time $T$
\begin{equation*}
    V(S, T) = f(S).
\end{equation*}

As in Black-Scholes model,
$\mu$ does not explicitly appear in the PDE,
its only present implicitly in the prices $S(t)$ of the assets.

To see that this is indeed a generalization of chapter \ref{ch:bs},
take $n = 1$ with a constant rate $r$, so that $S_0(t) = e^{rt}$, and $\sigma_1(S, t) = \sigma S_1$.
As $S_0$ is deterministic we can write the price as a function of time and the risky asset alone,
$f(t, S) = V(S_0(t), S, t)$, and the chain rule gives
\begin{equation*}
    \frac{\partial f}{\partial t}
        = \frac{\partial V}{\partial t} + r S_0 \frac{\partial V}{\partial S_0}.
\end{equation*}
Now \eqref{eq:replicating_portfolio_pde_2} gives the first term and \eqref{eq:replicating_portfolio_pde_3} the second,
\begin{equation*}
    \frac{\partial V}{\partial t} = - \frac{1}{2} \sigma^2 S^2 \frac{\partial^2 f}{\partial S^2},
    \qquad
    S_0 \frac{\partial V}{\partial S_0} = f(t, S) - S \frac{\partial f}{\partial S},
\end{equation*}
and substituting we recover the Black-Scholes PDE \eqref{eq:bs_pde}.
Note that the rate $r$ enters the pricing equation only through \eqref{eq:replicating_portfolio_pde_3},
i.e. through the position held in the riskless asset.

\section{Martingale method}\label{subsec:martingale_method}

The martingale method is an alternative framework to the PDE method,
that allows to price contingent claims in a more flexible way.

We will use \cite[Ch.~10]{bjork_arbitrage} as a reference for this section,
we start with a couple of definitions related with the concept of market and portfolio defined previously.
\begin{definition}[Attainable process]
    We say a price process $V(t)$ is \textbf{attainable} if there exists a self-financing portfolio $\theta(t)$ such that $V(t) = \theta(t)^\intercal S(t)$.
\end{definition}

The attainable process is a formalization of the dynamic replicating portfolio idea.
An other restriction on the portfolios we are going to consider is the following,
\begin{definition}[Admissible portfolio]
    A portfolio $\theta(t)$ is called \textbf{admissible} if exists $\alpha > 0$ such that
    \begin{equation*}
        \int^t_0 \theta(u)^\intercal dS(u) \geq -\alpha, \quad \forall t \in [0, T].
    \end{equation*}
\end{definition}
In real life we can interpret this condition as a limit on the amount the portfolio can lose or borrow,
as bank, clearing house or counterparty will not allow such risks and would force to cancel the deal.
Technically also allow to cancel out some other undesirable strategies that would prevent some key theorems to hold,
see \cite[Sec.~14.5 and 14.6]{steele_stochastic} for more details on such strategies.

Now we will formalize the idea of arbitrage in a market.
\begin{definition}[Arbitrage]\label{def:arbitrage}
    A self-financing and admissible portfolio $\theta(t)$ is an \textbf{arbitrage} if one of the following conditions hold for time $t$:
    \begin{enumerate}
        \item $\theta(0)^\intercal S(0) < 0$ and $\mathbb{P}(\theta(t)^\intercal S(t) \geq 0) = 1$,
        \item $\theta(0)^\intercal S(0) = 0$ and $\mathbb{P}(\theta(t)^\intercal S(t) \geq 0) = 1$ and $\mathbb{P}(\theta(t)^\intercal S(t) > 0) > 0$,
    \end{enumerate}
\end{definition}

Consider that the \textbf{risk-free} asset $S_0$ has a constant price of $1$, i.e. the risk-free rate is $r(S(t), t) \equiv 0$.
Under this assumption we introduce the \textbf{risk-neutral measure} $\Q$ as,
\begin{definition}\label{def:risk_neutral_measure}
    A probability measure $\Q$ on $\F_T$ is called a \textbf{risk-neutral measure} (also known as \textbf{equivalent martingale measure}) if
    \begin{itemize}
        \item $\Q$ is equivalent to $\mathbb{P}$ on $\F_T$, that is, both measures have the same null sets,
            \begin{equation*}
                \Q(A) = 0 \iff \mathbb{P}(A) = 0, \quad \forall A \in \F_T,
            \end{equation*}
            so they agree on which events are impossible, but not on how likely the possible ones are.
        \item All price process in the market $S_0, S_1, \dots, S_n$ are martingales under $\Q$ on the time interval $[0, T]$,
            \begin{equation*}
                \E^{\Q}[S_i(t) | \F_s] = S_i(s), \quad \forall\, 0 \leq s \leq t \leq T,
            \end{equation*}
            where $\E^{\Q}$ denotes the expectation taken under $\Q$.
    \end{itemize}
\end{definition}
We refer to $\mathbb{P}$ as the "real" probability measure and reflects the true probability of the price processes,
whereas $\Q$ is a probability measure under which the expected price of any asset is equal to the current one regardless of their risk (volatility).

Under this assumption we have the following theorem.
\begin{theorem}[First fundamental theorem of asset pricing]\label{thm:1_fundamental_thm}
    The market $\{ S_i(t) \}$ is arbitrage-free if and only if there exists an equivalent (local) martingale measure $\Q$.
\end{theorem}
We include the local martingales\footnote{
    A process $X$ is a \textbf{local martingale} if there is an increasing sequence of stopping times $\tau_m \to \infty$ such that every stopped process $X_{t \wedge \tau_m}$ is a martingale.
    Every martingale is a local martingale,
    but not the other way around,
    as the integrability condition may fail globally.
    See \cite[Ch.~7]{steele_stochastic}.
} as restricting to the martingales would eliminate too many processes,
that are of interest.

The restriction on the risk-free asset being constant can be dropped by considering the first price process a numeraire and normalizing the market.
For $S_0$ to be a numeraire we need that $S_0(t) > 0$ almost surely for all $t \geq 0$.
Using $S_0$ as numeraire we can consider the processes $Z_i = S_i(t) / S_0(t)$,
which are the \textbf{discounted} price processes of the assets.
\begin{definition}[Normalized market]\label{def:norm_market}
    A normalized market is a vector of  processes $Z$ given by
    \begin{equation*}
        Z_i(t) = \frac{S_i(t)}{S_0(t)}, \quad i = 1, \ldots, n.
    \end{equation*}
\end{definition}

The results presented so far do not depend on whether we work with the market $\{ S_i(t) \}$ as defined in \ref{def:market} or with the normalized market $\{ Z_i(t) \}$ of Definition \ref{def:norm_market}.
A portfolio is self-financing for one market if and only if it is self-financing for the other,
and the same holds for being an arbitrage,
see \cite[Ch.~10]{bjork_arbitrage}.
Theorem \ref{thm:1_fundamental_thm} therefore applies to any market admitting a numeraire, after normalizing it.

Notice that if the numeraire is the risk-free asset $S_0$ then the market $\{ S_i(t) \}$ is arbitrage-free if all assets $S_i(t)$ have a local rate of return $r(t)$ under $\Q$,
that is, if $\mu_i(S(t), t) = r(t) S_i(t)$ for all $i = 1, \ldots, n$ under $\Q$.
This is the sense in which $\Q$ is risk-neutral,
under it every asset drifts at the risk-free rate no matter its volatility.

Under the assumption of a market with no arbitrage,
a key question is whether a contingent claim can be replicated by a portfolio consisting of the market assets.
This yields the following theorem.
\begin{theorem}[Second fundamental theorem of asset pricing]\label{thm:2_fundamental_thm}
   Assuming absence of arbitrage,
   the market model is complete if and only if the martingale measure $\Q$ is unique.
\end{theorem}

Finally we can give a general formula to price contingent claims.
Under the assumption of a complete market,
the price at time $t$ of a contingent claim with payoff $f(S(T))$ is given by
\begin{equation}\label{eq:martingale_pricing}
    V(t, f(S(T))) = S_0(t)\, \E^{\Q} \left[ \left. \frac{f(S(T))}{S_0(T)} \,\right| \F_t \right],
\end{equation}
see \cite[Ch.~10]{bjork_arbitrage}.
This is just the martingale property of Definition \ref{def:risk_neutral_measure} read backwards,
the normalized price $V(t) / S_0(t)$ is a $\Q$-martingale whose terminal value is $f(S(T)) / S_0(T)$.
If the numeraire is the risk-free asset, $S_0(t) = e^{\int_0^t r(s) ds}$ and \eqref{eq:martingale_pricing} becomes
\begin{equation*}
    V(t, f(S(T))) = \E^{\Q} \left[ \left. e^{ - \int^T_t r(s) ds} f(S(T)) \,\right| \F_t \right].
\end{equation*}
